\documentclass[10pt,a4paper]{amsart}
\usepackage{amsmath,amssymb,amsthm,graphicx}
\usepackage{verbatim}
\newenvironment{reference}{\comment}{\endcomment}
%\newenvironment{reference}{}{}
\newenvironment{slogan}{\comment}{\endcomment}
\newenvironment{history}{\comment}{\endcomment}

% For commutative diagrams we use Xy-pic
\usepackage[all]{xy}

% We use 2cell for 2-commutative diagrams.
\xyoption{2cell}
% Theorem environments.
%
\theoremstyle{plain}
\newtheorem{theorem}[subsection]{Theorem}
\newtheorem{proposition}[subsection]{Proposition}
\newtheorem{lemma}[subsection]{Lemma}

\theoremstyle{definition}
\newtheorem{definition}[subsection]{Definition}
\newtheorem{example}[subsection]{Example}
\newtheorem{exercise}[subsection]{Exercise}
\newtheorem{situation}[subsection]{Situation}

\theoremstyle{remark}
\newtheorem{remark}[subsection]{Remark}
\newtheorem{remarks}[subsection]{Remarks}

\numberwithin{equation}{subsection}

% Macros
%
\def\lim{\mathop{\mathrm{lim}}\nolimits}
\def\colim{\mathop{\mathrm{colim}}\nolimits}
\def\Spec{\mathop{\mathrm{Spec}}}
\def\Hom{\mathop{\mathrm{Hom}}\nolimits}
\def\Ext{\mathop{\mathrm{Ext}}\nolimits}
\def\SheafHom{\mathop{\mathcal{H}\!\mathit{om}}\nolimits}
\def\SheafExt{\mathop{\mathcal{E}\!\mathit{xt}}\nolimits}
\def\Sch{\mathit{Sch}}
\def\Mor{\mathop{\mathrm{Mor}}\nolimits}
\def\Ob{\mathop{\mathrm{Ob}}\nolimits}
\def\Sh{\mathop{\mathit{Sh}}\nolimits}
\def\NL{\mathop{N\!L}\nolimits}
\def\CH{\mathop{\mathrm{CH}}\nolimits}
\def\proetale{{pro\text{-}\acute{e}tale}}
\def\etale{{\acute{e}tale}}
\def\QCoh{\mathit{QCoh}}
\def\Ker{\mathop{\mathrm{Ker}}}
\def\Im{\mathop{\mathrm{Im}}}
\def\Coker{\mathop{\mathrm{Coker}}}
\def\Coim{\mathop{\mathrm{Coim}}}

% Boxtimes
%
\DeclareMathSymbol{\boxtimes}{\mathbin}{AMSa}{"02}

%
% Macros for moduli stacks/spaces
%
\def\QCohstack{\mathcal{QC}\!\mathit{oh}}
\def\Cohstack{\mathcal{C}\!\mathit{oh}}
\def\Spacesstack{\mathcal{S}\!\mathit{paces}}
\def\Quotfunctor{\mathrm{Quot}}
\def\Hilbfunctor{\mathrm{Hilb}}
\def\Curvesstack{\mathcal{C}\!\mathit{urves}}
\def\Polarizedstack{\mathcal{P}\!\mathit{olarized}}
\def\Complexesstack{\mathcal{C}\!\mathit{omplexes}}
% \Pic is the operator that assigns to X its picard group, usage \Pic(X)
% \Picardstack_{X/B} denotes the Picard stack of X over B
% \Picardfunctor_{X/B} denotes the Picard functor of X over B
\def\Pic{\mathop{\mathrm{Pic}}\nolimits}
\def\Picardstack{\mathcal{P}\!\mathit{ic}}
\def\Picardfunctor{\mathrm{Pic}}
\def\Deformationcategory{\mathcal{D}\!\mathit{ef}}
\usepackage{cleveref}
\crefname{originallemma}{original lemma}{original lemmas}
\Crefname{originallemma}{Original lemma}{Original lemmas}
\makeatletter
\newlabel{lemma-adjoint}{{0778}{source}{}{sourcepointer-1}{}}
\newlabel{lemma-adjoint@cref}{{[originallemma][0][]0778}{[1][0][]source}{}{}{}}
\makeatother
\crefname{originalreference}{original reference}{original references}
\Crefname{originalreference}{Original reference}{Original references}
\makeatletter
\newlabel{stacks-sheaves-lemma-pullback-quasi-coherent}{{06WH}{source}{}{sourcepointer-2}{}}
\newlabel{stacks-sheaves-lemma-pullback-quasi-coherent@cref}{{[originalreference][0][]06WH}{[1][0][]source}{}{}{}}
\makeatother
\crefname{originalproposition}{original proposition}{original propositions}
\Crefname{originalproposition}{Original proposition}{Original propositions}
\makeatletter
\newlabel{proposition-loc-qcoh-flat-base-change}{{0771}{source}{}{sourcepointer-3}{}}
\newlabel{proposition-loc-qcoh-flat-base-change@cref}{{[originalproposition][0][]0771}{[1][0][]source}{}{}{}}
\makeatother
\crefname{originallemma}{original lemma}{original lemmas}
\Crefname{originallemma}{Original lemma}{Original lemmas}
\makeatletter
\newlabel{lemma-adjoint-kernel-parasitic}{{0779}{source}{}{sourcepointer-4}{}}
\newlabel{lemma-adjoint-kernel-parasitic@cref}{{[originallemma][0][]0779}{[1][0][]source}{}{}{}}
\makeatother
\crefname{originallemma}{original lemma}{original lemmas}
\Crefname{originallemma}{Original lemma}{Original lemmas}
\makeatletter
\newlabel{lemma-parasitic}{{0774}{source}{}{sourcepointer-5}{}}
\newlabel{lemma-parasitic@cref}{{[originallemma][0][]0774}{[1][0][]source}{}{}{}}
\makeatother
\crefname{originallemma}{original lemma}{original lemmas}
\Crefname{originallemma}{Original lemma}{Original lemmas}
\makeatletter
\newlabel{lemma-exact-sequence-quasi-coherent-parasitic-cohomology}{{0776}{source}{}{sourcepointer-6}{}}
\newlabel{lemma-exact-sequence-quasi-coherent-parasitic-cohomology@cref}{{[originallemma][0][]0776}{[1][0][]source}{}{}{}}
\makeatother
\title{Flat pullback commutes with the Q functor}
\author{The Stacks Project Authors}
\date{stacks-project:stacks-cohomology.tex:local-snapshot-20261012; GNU FDL 1.2 or later}
\begin{document}
\maketitle
\noindent\textit{Editorial note.} All delivered mathematical text is the author's own, copied verbatim from the pinned Stacks Project chapter file stacks-cohomology.tex; no statement or proof step was rewritten, reordered or removed.\par\medskip
\noindent\textit{Editorial note.} The counted span cites the lemmas recorded in assembly.external\_references with canonical tags as navigation pointers.\par\medskip
\noindent\textit{Allowed standard prerequisites (editorial).}\begin{itemize}
\item Parasitic kernel and cokernel of the adjunction map. Cohomology of Algebraic Stacks, Lemma stacks-cohomology-lemma-adjoint-kernel-parasitic (2). Identifies the comparison map's defect.
\item Flat pullback preserves parasitic modules. Cohomology of Algebraic Stacks, Lemma stacks-cohomology-lemma-parasitic (1). Keeps the defect parasitic after pullback.
\item Parasitic cohomology exactness criterion. Cohomology of Algebraic Stacks, Lemma stacks-cohomology-lemma-exact-sequence-quasi-coherent-parasitic-cohomology. Parasitic kernel and cokernel in QCoh imply the map is an isomorphism.
\end{itemize}\medskip
\ifdefined\hypertarget\hypertarget{sourcepointer-1}{}\fi
\noindent\textit{Source reference.} Original lemma 0778: sources/stacks-chow/tags.txt, source lines 9193--9193. Cross-chapter standard prerequisite cited by the author inside the counted span; it is used as a pointer and its own proof is not part of this item.\par\medskip
\ifdefined\hypertarget\hypertarget{sourcepointer-2}{}\fi
\noindent\textit{Source reference.} Original reference 06WH: sources/stacks-chow/tags.txt, source lines 8816--8816. Cross-chapter standard prerequisite cited by the author inside the counted span; it is used as a pointer and its own proof is not part of this item.\par\medskip
\ifdefined\hypertarget\hypertarget{sourcepointer-3}{}\fi
\noindent\textit{Source reference.} Original proposition 0771: sources/stacks-chow/tags.txt, source lines 9186--9186. Cross-chapter standard prerequisite cited by the author inside the counted span; it is used as a pointer and its own proof is not part of this item.\par\medskip
\ifdefined\hypertarget\hypertarget{sourcepointer-4}{}\fi
\noindent\textit{Source reference.} Original lemma 0779: sources/stacks-chow/tags.txt, source lines 9194--9194. Cross-chapter standard prerequisite cited by the author inside the counted span; it is used as a pointer and its own proof is not part of this item.\par\medskip
\ifdefined\hypertarget\hypertarget{sourcepointer-5}{}\fi
\noindent\textit{Source reference.} Original lemma 0774: sources/stacks-chow/tags.txt, source lines 9189--9189. Cross-chapter standard prerequisite cited by the author inside the counted span; it is used as a pointer and its own proof is not part of this item.\par\medskip
\ifdefined\hypertarget\hypertarget{sourcepointer-6}{}\fi
\noindent\textit{Source reference.} Original lemma 0776: sources/stacks-chow/tags.txt, source lines 9191--9191. Cross-chapter standard prerequisite cited by the author inside the counted span; it is used as a pointer and its own proof is not part of this item.\par\medskip
\begin{lemma}
\label{lemma-flat-morphism-and-Q}
Let $f : \mathcal{X} \to \mathcal{Y}$ be a flat morphism of algebraic stacks.
Then $Q_\mathcal{X} \circ f^* = f^* \circ Q_\mathcal{Y}$ where
$Q_\mathcal{X}$ and $Q_\mathcal{Y}$ are as in Lemma \ref{lemma-adjoint}.
\end{lemma}

\begin{proof}
Observe that $f^*$ preserves both $\QCoh$ and $\textit{LQCoh}^{fbc}$,
see Sheaves on Stacks, Lemma
\ref{stacks-sheaves-lemma-pullback-quasi-coherent} and
Proposition \ref{proposition-loc-qcoh-flat-base-change}.
If $\mathcal{F}$ is in $\textit{LQCoh}^{fbc}(\mathcal{O}_\mathcal{Y})$
then $Q_\mathcal{Y}(\mathcal{F}) \to \mathcal{F}$
has parasitic kernel and cokernel
by Lemma \ref{lemma-adjoint-kernel-parasitic}.
As $f$ is flat we get that $f^*Q_\mathcal{Y}(\mathcal{F}) \to f^*\mathcal{F}$
has parasitic kernel and cokernel by
Lemma \ref{lemma-parasitic}.
Thus the induced map
$f^*Q_\mathcal{Y}(\mathcal{F}) \to Q_\mathcal{X}(f^*\mathcal{F})$
has parasitic kernel and cokernel and hence is an isomorphism
for example by
Lemma \ref{lemma-exact-sequence-quasi-coherent-parasitic-cohomology}.
\end{proof}

\end{document}
